\section{The full Guralnick--Saxl strengthening}
\label{sec:scalar-unramified-radical}

\begin{corollary}[D-BI-14a: a quadratic bound in the ramified prime]
Under the primitive, tensor-indecomposable hypotheses of D-BI-14,
put $d=p-1$.  Then
\[
                       e\ge\frac{d(25d-22)}4.              \tag{BI.50a}
\]
For $p=7,11,19,43,67,163$ these bounds are respectively
\[
              192,\ 570,\ 1926,\ 10794,\ 26862,\ 163134.
\]
\end{corollary}
\begin{proof}
We use the complete statement of Guralnick--Saxl
\cite[Theorem~7.1 and Remark~(3), pp.~548--549;
Lemma~7.4, pp.~550--551]{BI:GS}, rather
than just its bireflection consequence.  Their invariant
$\nu_J(V)$ is the minimum moving rank of $\lambda g\ne1$,
where $g\in J$ and $\lambda$ is an arbitrary scalar.
Outside the listed exceptions their theorem gives
\[
                    \nu_J(V)>\max\{2,\sqrt e/2\}.
\]
In characteristic zero and degree $e\ge42$, all the exceptional
families except the scalar deleted alternating/symmetric family
are absent: the natural finite classical cases are in defining
positive characteristic, and all other listed exceptions have
degree at most ten.  D-BI-08 excludes the scalar deleted family
also for a larger inducing field.  The theorem's tensor convention
is exactly the convention in D-BI-14; it includes a group that
permutes tensor factors, and the extra tensor-cube exception has
degree eight.  Thus every nonidentity element of $J$ has moving
rank at least
\[
                        q=\lfloor\sqrt e/2\rfloor+1.
\]
Apply (BI.48) to obtain $e\ge(5q-3)d$.  The definition of $q$
also gives $4(q-1)^2\le e<4q^2$, whence
\[
                         4q^2-5dq+3d>0.                  \tag{BI.50b}
\]
Here $d\ge6$ and $d\equiv2\pmod4$.  Put
$q_0=(5d-2)/4$, an integer.  The upward quadratic in (BI.50b)
is negative at $q=1$ and at $q=q_0-1$:
\[
 4-2d<0,\qquad
 4(q_0-1)^2-5d(q_0-1)+3d=9-\tfrac92d<0.
\]
By convexity it is negative throughout $1\le q\le q_0-1$.
Therefore $q\ge q_0$, and
\[
           e\ge(5q_0-3)d=\frac{d(25d-22)}4,
\]
as asserted.  This is only a necessary degree bound: no existence
or dominance follows at the numerical threshold.
\end{proof}
